
\chapter{Métodos de Inferência Probabilística para Reconstrução Filogenômica}

\section{Introdução}

A reconstrução de relações filogenéticas a partir de dados genômicos situa-se na intersecção da biologia evolutiva, estatística e ciência computacional. Nos últimos três decênios, o paradigma analítico dominante em filogenética deslocou-se de métodos baseados em parcimônia para abordagens probabilísticas --- máxima verossimilhança (MV) e inferência Bayesiana (IB). Este deslocamento foi impulsionado pela explosão de dados de sequências moleculares, pelo desenvolvimento de modelos cada vez mais sofisticados de evolução de sequências e pela disponibilidade de poder computacional que teria parecido fantástico aos sistematas dos anos 1990.

Como Wheeler (2012) nos lembra, ``não existe uma única resposta para `qual é o melhor método?' da mesma forma que não existe uma única resposta para `qual é a melhor árvore?' Um critério deve ser especificado em ambos os casos'' \parencite{wheeler2012}. A escolha de um marco probabilístico é ela própria uma hipótese sobre como a inferência deve proceder --- e é uma que tem sido vigorosamente contestada em fundamentos teóricos, filosóficos e práticos.

\section{Fundamentos Teóricos da Probabilidade em Filogenética}

\subsection{A Função de Verossimilhança}

A verossimilhança de uma árvore $T$ dado dados $D$ é proporcional à probabilidade dos dados dada a árvore \parencite{edwards1972}:

$$l(T|D) \propto \Pr(D|T)$$

Um método sistemático de MV seleciona a árvore $T$ que maximiza $\Pr(D|T)$. Esta maximização requer especificar um conjunto de \textit{parâmetros de incômodo} (\textit{nuisance parameters}) $\theta$, definidos como todas as quantidades necessárias para calcular $\Pr(D|T)$ além dos dados e da topologia da árvore \parencite{wheeler2012}. Os três parâmetros de incômodo mais consequentes são: (1) o modelo de transformação de caracteres; (2) parâmetros de ramo (comprimentos de ramo); e (3) a distribuição de taxas de mudança entre caracteres.

\subsection{O Modelo Markov Geral}

A descrição reversível mais geral de mudança de caráter para $n$ estados é especificada por uma matriz de taxa instantânea $R$ e um vetor de frequência estacionária $\symbf{\pi}$. As condições de simetria impostas em $R$ são:

$$\forall i,\; R_{ii} = 0; \quad \forall i,j,\; R_{ij} = R_{ji}; \quad \sum_{i=1}^{n}\sum_{j=1}^{n} \pi_i \cdot \pi_j \cdot R_{ij} = 1$$

\textbf{É assim que leríamos essa equação, passo-a-passo:}

\begin{itemize}
    \item \textbf{Parte 1:} ``Para todo $i$, erre-sub-ii é igual a zero;''
    \item \textbf{Parte 2:} ``Para todo $i$ e $j$, erre-sub-ij é igual a erre-sub-ji;''
    \item \textbf{Parte 3:} ``O somatório duplo, com $i$ e $j$ de 1 até $n$, de pi-sub-i vezes pi-sub-j vezes erre-sub-ij é igual a 1.''
\end{itemize}

\textbf{Veja algumas observações de pronúncia e contexto:}

\begin{center}
\renewcommand{\arraystretch}{1.5}
\begin{tabular}{lp{10cm}}
\toprule
\textbf{Símbolo} & \textbf{Descrição e Pronúncia} \\
\midrule
$\forall$ & Lido como ``Para todo'' ou ``Qualquer que seja''. \\
$R_{ii}, R_{ij}$ & Pronuncia-se a letra e os índices (``erre-ii'' ou ``erre-ij''). \\
$\pi_i$ & O som de $\pi$ em português é sempre ``pí'', nunca ``pái''. \\
$\sum \sum$ & Chamado de ``Somatório duplo''. \\
$n$ & Lido simplesmente como a letra ``ene''. \\
\bottomrule
\end{tabular}
\end{center}

Estas matrizes são combinadas na matriz de taxa $Q$:

$$Q_{ij} = \begin{cases} R_{i,j} \cdot \pi_j & \text{se } i \neq j \\ -\sum_{m=1}^{n} R_{i,m} \cdot \pi_m & \text{se } i = j \end{cases}$$

A probabilidade de mudança entre estados $i$ e $j$ em tempo $t$ é:

$$P_{i,j}(t) = \sum_{m=1}^{n} e^{\lambda_m t} \cdot U_{m,i} \cdot U^{-1}_{j,m}$$

onde $\lambda_m$ (lâmbda-sub-$m$) são os autovalores de $Q$, $U$ é a matriz de autovetores e $U^{-1}$ sua inversa.

\section{Classificação de Métodos de Verossimilhança Máxima}

\subsection{O Conceito de Parâmetros de Incômodo}

Wheeler (2012) apresenta uma classificação hierárquica dos métodos de verossimilhança usados em sistemática. A ideia-chave é que cada método difere em \textit{como trata os parâmetros de incômodo e as atribuições de estados ancestrais}. Este é um ponto fundamental frequentemente negligenciado: diferentes formulações do que exatamente está sendo maximizado --- médias sobre estados (MVM), melhores estados em nós (MVP), ou melhores caminhos evolutivos completos (VPC) --- levam a métodos genuinamente diferentes com propriedades distintas.

\subsection{Máxima Verossimilhança Integrada (MVI)}

Se conhecêssemos a distribuição de parâmetros de incômodo, $\Phi(\theta|T)$, poderíamos integrá-los fora do espaço de parâmetros $\Theta$:

$$p(D|T) = \int p(D|T, \theta) d\Phi(\theta|T)$$

\textbf{Assim leríamos esta equação, passo-a-passo:}

\begin{itemize}
    \item \textbf{Lado Esquerdo:} ``A probabilidade pê ($p$) de dê ($D$) dado tê ($T$) é igual a...''
    \item \textbf{Lado Direito:} ``A integral da probabilidade pê ($p$) de dê ($D$) dado tê ($T$) e teta ($\theta$), em relação à medida fi ($\Phi$) de teta ($\theta$) dado tê ($T$).''
    \item \textbf{Leitura Fluida:} ``A probabilidade de D dado T é a integral da verossimilhança de D dado T e teta ($\theta$), integrada sobre a distribuição \textit{a priori} (ou medida) de teta ($\theta$) dado T.''
\end{itemize}

\vspace{0.5cm}

\textbf{Algumas observações de pronúncia e contexto:}

\begin{center}
\renewcommand{\arraystretch}{1.5}
\begin{tabular}{lp{10cm}}
\hline
\textbf{Símbolo} & \textbf{Descrição e Pronúncia} \\ \hline
$p(A|B)$ & O traço vertical ``$|$'' é sempre lido como \textbf{``dado''}. \\
$\int$ & Lido como \textbf{``Integral''}. \\
$\theta$ & Letra grega \textbf{``teta''}. \\
$\Phi$ & Letra grega \textbf{``fi''} (maiúsculo). \\
$d\Phi$ & O ``$d$'' indica o diferencial, lido como \textbf{``em relação a''} ou \textbf{``na medida de''}. \\ \hline
\end{tabular}
\end{center}

A árvore que maximiza esta probabilidade integrada é a árvore de \textbf{Máxima Verossimilhança Integrada} (MVI). Esta é a abordagem teoricamente mais satisfatória porque conta para nossa incerteza nos parâmetros de incômodo ao invés de escolher um único conjunto ``melhor''.

\subsection{Máxima A Posteriori (MAP)}

Se também temos uma \textbf{distribuição priori sobre árvores}, $p(T)$, podemos perguntar: qual árvore tem a maior probabilidade posterior? Székelyi e Steel (1999) mostraram que o método com a maior expectativa de retornar a árvore verdadeira seleciona a árvore maximizando $p(T) \times \Pr(D|T)$. Esta é a árvore de \textbf{Máxima A Posteriori} (MAP) --- a estimativa Bayesiana.

Aqui está a conexão importante: quando todas as prioris de árvore são \textbf{iguais} (uma priori flat, não-informativa), a árvore MAP é idêntica à árvore MVI. Como Wheeler nota, ``o uso de prioris de árvore não-uniformes (tais como empíricas ou Yule) quebra esta identidade'' \parencite{wheeler2012}.

\subsection{Máxima Verossimilhança Relativa (MVR) e Suas Três Variantes}

Na prática, raramente conhecemos a distribuição de parâmetros de incômodo. Ao invés de integrá-los, \textbf{otimizamos}: escolhemos $\theta$ para maximizar $p(D|T, \theta)$. Isto é chamado \textbf{Máxima Verossimilhança Relativa} (MVR), e é ``em geral, a metodologia usada em análises empíricas'' \parencite{wheeler2012}.

\subsubsection{Máxima Verossimilhança Média (MVM)}

Este é o abordagem mais comumente usada. Ela \textbf{soma sobre todas as atribuições de estados possíveis de vértice}, ponderadas por suas probabilidades. Quando as pessoas referem-se a ``MV'' em análises filogenéticas típicas, geralmente significam MVM --- e Wheeler nota que ``para o restante desta discussão, quando falarmos de métodos de MV, estaremos referindo-nos a MVM'' \parencite{wheeler2012}.

Em MVM, os estados ancestrais são parâmetros de incômodo. Não nos importamos realmente qual estado um ancestral tinha; nos importamos com a verossimilhança da \textit{árvore e dos dados}. A forma estatisticamente correta de remover quantidades desconhecidas de um cálculo de probabilidade é \textbf{marginalizar} (isto é, somar ou integrar) sobre elas, não otimizá-las.

\subsubsection{Máxima Verossimilhança Mais Parcimoniosa (MVP)}

Também sugerida por Barry e Hartigan (1987), MVP \textbf{atribui estados de vértice específicos} (junto com outros parâmetros) tais que a verossimilhança geral da árvore é maximizada. O nome insinua uma conexão com parcimônia, e de fato a ideia de selecionar estados ancestrais ``melhores'' soa similar. Porém, há um detalhe: em MVP, as \textbf{probabilidades de ramo são as mesmas sobre todos os caracteres}, então ``MVP não escolherá (em geral) a mesma árvore que parcimônia'' \parencite{wheeler2012}.

MVP difere de MVM em que otimiza sobre estados ancestrais. Ela escolhe a atribuição única melhor. Em estatística, substituir uma integral por uma otimização é conhecido como tomar um estimador ``plug-in'' ou de ponto, e tende a \textbf{superestimar a evidência} para uma árvore particular.

\subsubsection{Verossimilhança de Caminho Evolutivo (VCE)}

Originalmente proposta por Farris (1973), VCE vai além de MVP. Ao invés de apenas especificar estados de vértice ancestrais, VCE especifica \textbf{a sequência inteira de estados de caráter intermediários entre vértices} tais que a verossimilhança geral da árvore é maximizada. Pense nela como maximizar a verossimilhança ao longo de cada único ``caminho'' evolutivo na árvore.

Um resultado notável: a árvore que maximiza VCE \textbf{é precisamente a árvore mais parcimoniosa}. Wheeler afirma isto diretamente:

\begin{quote}
``A árvore que maximiza esta forma de verossimilhança é precisamente a árvore mais parcimoniosa. Este resultado se sustenta para um conjunto amplo e robusto de assunções (não há exigência de taxas baixas ou homogêneas de mudança de caráter, por exemplo)'' \parencite{wheeler2012}
\end{quote}

\subsection{O Paradoxo: Felsenstein vs. Farris}

A relação entre parcimônia e verossimilhança cria o que parece ser uma contradição. Felsenstein (1973, 1978) argumentou que métodos de MV são estatisticamente consistentes e mostrou que a parcimônia pode ser inconsistente sob certas condições (a ``Zona de Felsenstein,'' onde ramos longos se atraem mutuamente). Mas Farris (1973a) mostrou que parcimônia \textit{é} um método de MV sob a formulação VCE. Se MV é consistente e parcimônia é MV, então parcimônia deveria ser consistente também---mas Felsenstein mostrou que nem sempre é. Wheeler (2012, p.~218) resolve este aparente paradoxo reconhecendo que Farris e Felsenstein estavam discutindo \textbf{formas diferentes de verossimilhança}: a MV de Felsenstein refere-se a MVM (média sobre estados ancestrais com um modelo comum entre caracteres), enquanto a MV de Farris refere-se a VCE (maximização sobre o caminho inteiro de mudanças evolutivas). São problemas de otimização genuinamente diferentes que podem produzir árvores ótimas diferentes. Os resultados de consistência de Felsenstein aplicam-se a MVM, não a VCE, e assim a contradição se dissolve.

\subsection{Conexões Adicionais: Goldman (1990) e Tuffley \& Steel (1997)}

A ligação entre parcimônia e verossimilhança foi explorada por Goldman (1990), que mostrou que quando os pesos de aresta são constantes sobre a árvore, verossimilhança e parcimônia convergem \parencite{goldman1990}. Tuffley e Steel (1997) introduziram o \textbf{modelo Sem-Mecanismo-Comum} (\textit{No-Common-Mechanism}, NCM), onde cada caráter tem uma taxa potencialmente única que pode variar entre arestas \parencite{tuffley1997}. Sob este modelo e uma estrutura tipo Neyman, a verossimilhança global torna-se uma função direta do comprimento de parcimônia. Porém, como Wheeler (2012, p.~224) alerta, esta equivalência se sustenta estritamente apenas ``quando o número de estados de cada caráter ($r_i$) é uma constante sobre o conjunto de dados.'' Quando $r_i$ varia entre caracteres, a verossimilhança NCM introduz um fator de ponderação específico por caráter, de modo que ``NCM pode ser visto como um esquema de ponderação de caracteres baseado em verossimilhança em análises de parcimônia'' \parencite{wheeler2012}.

É importante distinguir VCE do modelo NCM: VCE é uma variante em \textit{como estados ancestrais são tratados}, situando-se dentro da classificação de Wheeler como um sabor de Máxima Verossimilhança Relativa. O modelo NCM, em contraste, é uma variante no \textit{modelo evolutivo em si}. VCE produz uma equivalência exata com parcimônia sob condições muito gerais (sem exigência de taxas baixas ou homogêneas). NCM produz equivalência com parcimônia apenas condicionalmente---quando $r_i$ é constante entre caracteres. Apesar dessas diferenças, ambos os resultados convergem na mesma conclusão: que a parcimônia tem uma fundação principiada dentro da estrutura de verossimilhança.

\section{A Perspectiva Bayesiana}

\subsection{Teorema de Bayes e Inferência Posterior}

A inferência Bayesiana em filogenética reformula inferência de árvore usando o Teorema de Bayes. A probabilidade posterior de uma topologia $T$ dado dados $D$, parâmetros de substituição $\theta$, e pesos de ramo $t_T$ é:

$$p(T|D, \theta, t_T) = \frac{p(T) \cdot p(D|T, \theta, t_T)}{\sum_{T_j \in \tau} p(T_j) \cdot p(D|T_j, \theta, t_T)}$$

\textbf{Assim leríamos esta equação, passo-a-passo:}

\begin{itemize}
    \item \textbf{Lado Esquerdo:} ``A probabilidade pê ($p$) de tê ($T$) dado dê ($D$), teta ($\theta$) e tê sub-tê ($t_T$) é igual a...''
    \item \textbf{Numerador:} ``A probabilidade \textit{a priori} de tê ($T$) vezes a verossimilhança de dê ($D$) dado tê ($T$), teta ($\theta$) e tê sub-tê ($t_T$)...''
    \item \textbf{Denominador:} ``...dividido pelo somatório, para toda topologia tê-j ($T_j$) pertencente ao conjunto tau ($\tau$), da probabilidade de tê-sub-j ($T_j$) vezes a verossimilhança de dê ($D$) dado tê-sub-j ($T_j$), teta ($\theta$) e tê sub-tê ($t_T$).''
    \item \textbf{Leitura Fluida:} ``A posterior da árvore $T$ é o produto da sua priori pela sua verossimilhança, normalizado pelo somatório das probabilidades de todas as topologias possíveis no espaço das árvores $\tau$.''
\end{itemize}

\vspace{0.5cm}

\textbf{Algumas observações de pronúncia e contexto:}

\begin{center}
\renewcommand{\arraystretch}{1.5}
\begin{tabular}{lp{10cm}}
\hline
\textbf{Símbolo} & \textbf{Descrição e Pronúncia} \\ \hline
$p(T)$ & Frequentemente chamada de \textbf{``priori''} da topologia. \\
$t_T$ & Lido como \textbf{``tê sub-tê''}, representa os comprimentos dos ramos da árvore $T$. \\
$\sum_{T_j \in \tau}$ & Lido como \textbf{``somatório para cada tê-j pertencente a tau''}. \\
$\tau$ & Letra grega \textbf{``tau''}, aqui representando o conjunto de todas as topologias possíveis. \\
$\frac{A}{B}$ & Em voz alta, dizemos o numerador e usamos a expressão \textbf{``sobre''} ou \textbf{``dividido por''} para o denominador (neste caso, $A$ sobre $B$ ou $A$ dividido por $B$). \\ \hline
\end{tabular}
\end{center}

Integração sobre parâmetros de incômodo remove a condicionalidade sobre eles:

$$p(T|D) = \frac{\iint p(\theta), p(T|\theta), p(t_T|\theta, T), p(D|\theta, T, t_T)\; dt_T, d\theta}{\sum_{T_j \in \tau} \iint p(\theta), p(T_j|\theta), p(t_{T_j}|\theta, T_j), p(D|\theta, T_j, t_{T_j})\; dt_{T_j}, d\theta}$$

\textbf{Assim leríamos esta equação, passo-a-passo:}

\begin{itemize}
    \item \textbf{Lado Esquerdo:} ``A probabilidade posterior da árvore ($T$) dados apenas os dados ($D$).''
    \item \textbf{Numerador:} ``A integral dupla sobre os parâmetros teta ($\theta$) e pesos de ramo ($t_T$) da probabilidade conjunta de todos os componentes do modelo.''
    \item \textbf{Denominador:} ``O somatório, para todas as árvores possíveis ($T_j$), da mesma integral dupla realizada no numerador.''
    \item \textbf{Leitura Fluida:} ``A probabilidade da topologia $T$ é a sua verossimilhança marginalizada sobre todos os possíveis valores de parâmetros e comprimentos de ramo, normalizada pela soma das verossimilhanças marginais de todas as outras topologias.''
\end{itemize}

\vspace{0.5cm}

\textbf{Algumas observações de pronúncia e contexto:}

\begin{center}
\renewcommand{\arraystretch}{1.5}
\begin{tabular}{lp{10cm}}
\hline
\textbf{Símbolo} & \textbf{Descrição e Pronúncia} \\ \hline
$\iint$ & Lido como \textbf{``Integral dupla''}. Representa a integração sobre múltiplos parâmetros simultaneamente. \\
$dt_T, d\theta$ & Indicam as variáveis de integração. Lemos como \textbf{``em relação a tê-sub-tê e teta''}. \\
Marginalização & O processo de usar integrais para ``remover'' parâmetros de incômodo ($\theta, t_T$) e focar apenas no que interessa ($T$). \\
$p(T|D)$ & Note que $\theta$ e $t_T$ sumiram do lado esquerdo; isso ocorre porque foram \textbf{``integrados''} para fora da equação. \\ \hline
\end{tabular}
\end{center}

\subsection{Convergência Bayesiana e o Teorema de Bernstein-von Mises}

Uma garantia teórica fundamental frequentemente citada em análise Bayesiana é que, dados dados suficientes, a distribuição posterior convergerá à ``verdade''. Isto é formalizado pelo \textbf{Teorema de Bernstein-von Mises} (BvM). O Teorema BvM afirma que sob condições de regularidade específicas, conforme a quantidade de dados (em filogenética, o comprimento de sequência $n$) aproxima-se do infinito, a distribuição posterior converge para uma distribuição normal multivariada centrada nos valores de parâmetros verdadeiros (ou no estimador de máxima verossimilhança), com variância encolhendo a zero \parencite{vandervaart1998}.

Porém, é crítico reconhecer que esta garantia de convergência é tanto \textbf{assintótica} quanto \textbf{condicional}:

\begin{itemize}
    \item \textbf{Assintótica.} Aplica-se estritamente apenas no limite conforme $n \to \infty$. Para qualquer alinhamento finito (que é tudo que temos), a garantia é uma aproximação cuja qualidade depende de quão grande $n$ é relativo à complexidade do problema.
    \item \textbf{Condicional.} A formulação padrão do Teorema BvM requer que o mecanismo gerador de dados verdadeiro (i.e., o processo biológico real que produziu as sequências) situe-se perfeitamente dentro do espaço de modelos predefinidos sendo explorado. Se o modelo verdadeiro $M_{\text{verdadeiro}}$ não é um elemento do espaço de modelos $\mathcal{M}$, as garantias fundamentais do Teorema BvM padrão não se aplicam diretamente.
\end{itemize}

\subsection{Topologias Discretas e Espaço BHV}

Uma sutileza adicional frequentemente negligenciada: o Teorema BvM clássico aplica-se a parâmetros contínuos em espaço euclidiano. Topologias de árvore, porém, são \textbf{objetos combinatórios discretos}. O espaço de parâmetros filogenético não é uma variedade suave. É, ao invés, uma coleção de ortantes (um por topologia) colados ao longo de suas fronteiras, formando o espaço de árvore Billera-Holmes-Vogtmann (BHV) \parencite{billera2001}.

O Teorema BvM pode ser aplicado aos parâmetros contínuos (comprimentos de ramo, parâmetros do modelo de substituição) \textit{condicional em uma topologia fixada}, mas a convergência de probabilidade posterior sobre topologias é um problema separado e mais difícil. Para a topologia discreta em si, o comportamento de convergência depende de condições adicionais, incluindo se a topologia verdadeira é identificável sob o modelo assumido, e se os dados fornecem sinal suficiente para distinguir entre topologias concorrentes.

\begin{tcolorbox}[
    colback=green!5!white, 
    colframe=green!50!black, 
    title=\textbf{Explicando de forma simples: O Problema das Árvores},
    arc=2mm, % Arredondamento dos cantos
    boxrule=0.5pt
]
Imagine que você está tentando mapear o caminho de uma estrada. Em uma geometria comum, você pode se mover suavemente para qualquer lado. Mas, na filogenética, mudar a estrutura de uma árvore (a "topologia") é como pular de um universo para outro. 

Cada árvore possível é como uma sala separada. Você pode caminhar livremente dentro da sala (ajustando o tamanho dos ramos), mas para mudar quem é "parente" de quem, você precisa atravessar uma porta estreita que conecta essas salas. O \textbf{Espaço BHV} é o nome matemático para esse "labirinto de salas" coladas umas às outras. O desafio é que as regras normais de probabilidade funcionam bem dentro de cada sala, mas ficam muito mais complicadas quando tentamos decidir em qual sala os dados realmente deveriam estar.
\end{tcolorbox}

\begin{tcolorbox}[
    colback=red!5!white, 
    colframe=red!75!black, 
    title=\textbf{Consequências Práticas para a Pesquisa},
    arc=2mm
]
\begin{itemize}
    \item \textbf{O Problema da Convergência (MCMC):} Como o espaço das árvores é um ``labirinto de salas'', os algoritmos (como os do MrBayes ou BEAST) muitas vezes ficam ``presos'' em uma sala (uma topologia) que parece boa, mas não conseguem encontrar a porta para uma sala ainda melhor. É por isso que exigimos milhões de gerações e o uso de múltiplas cadeias frias e quentes (MCMCMC).
    \item \textbf{Incerteza não-Euclidiana:} Você não pode simplesmente tirar a ``média'' de duas árvores diferentes para obter uma árvore média. Na prática, usamos o \textit{Consenso de Maioria} ou \textit{Árvores de Máxima Credibilidade de Clado} porque a média aritmética não faz sentido em espaços discretos.
    \item \textbf{Sinal Filogenético vs. Ruído:} Se a ``verdade'' estiver exatamente na fronteira entre duas salas (ex: uma radiação muito rápida onde o ramo interno é quase zero), o modelo pode oscilar infinitamente entre topologias, gerando suporte baixo (valores de posterior baixos) não por erro, mas pela própria geometria do espaço.
\end{itemize}
\end{tcolorbox}

\section{Especificação Incorreta de Modelo}

\subsection{Convergência para a Melhor Resposta Errada}

Em filogenética, podemos afirmar com confiança que o modelo é \textit{sempre} especificado incorretamente em algum grau. Evolução de sequência biológica é um processo profundamente complexo impulsionado por taxas de mutação heterogêneas, seleção natural operando em múltiplos níveis, assortimento incompleto de linhagens, transferência horizontal de genes, dinâmica populacional mutável, e processos de substituição dependentes de contexto. Modelos Markov padrão de evolução de sequência (\textit{e.g.}, GTR+$\Gamma$) são abstrações matematicamente tratáveis.

Quando o modelo é especificado incorretamente ($M_{\text{verdadeiro}} \notin \mathcal{M}$), devemos perguntar: a que a distribuição posterior Bayesiana converge conforme o comprimento de sequência aumenta?

Sob especificação incorreta de modelo, a distribuição posterior não converge para a verdade absoluta. Ao invés, conforme $n \to \infty$, a probabilidade posterior concentra-se nos valores de parâmetros dentro do espaço de modelos explorado que minimizam a divergência de Kullback-Leibler (KL) da distribuição verdadeira geradora de dados. Formalmente, convergência é para:

$$M^* = \arg\min_{M \in \mathcal{M}} D_{KL}(P_{\text{verdadeiro}} | P_M)$$

\textbf{Assim leríamos esta equação, passo-a-passo:}

\begin{itemize}
    \item \textbf{Lado Esquerdo:} ``M-estrela ($M^*$) é igual ao argumento do mínimo...''
    \item \textbf{Lado Direito:} ``...sobre o conjunto de modelos M ($\mathcal{M}$), da Divergência de Kullback-Leibler ($D_{KL}$) entre a distribuição verdadeira ($P_{\text{verdadeiro}}$) e a distribuição do modelo ($P_M$).''
    \item \textbf{Leitura Fluida:} ``O modelo ideal $M^*$ é aquele, dentro do nosso conjunto de modelos possíveis, que minimiza a perda de informação (Divergência KL) em relação à realidade dos dados.''
\end{itemize}

\vspace{0.5cm}

\textbf{Algumas observações de pronúncia e contexto:}

\begin{center}
\renewcommand{\arraystretch}{1.5}
\begin{tabular}{lp{10cm}}
\hline
\textbf{Símbolo} & \textbf{Descrição e Pronúncia} \\ \hline
$M^*$ & O asterisco indica o valor \textbf{ótimo} ou a solução do problema. \\
$\arg\min$ & Não queremos o valor mínimo da distância, mas sim \textbf{qual modelo} gera esse mínimo. \\
$\mathcal{M}$ & Letra caligráfica representando o \textbf{espaço de busca} (todos os modelos que você está testando). \\
$D_{KL}(P|Q)$ & A \textbf{Divergência de Kullback-Leibler}. Mede a ``distância'' (ou surpresa extra) entre duas distribuições. \\ \hline
\end{tabular}
\end{center}

Este $M^*$ é às vezes chamado o valor de parâmetro ``pseudo-verdadeiro'': a melhor aproximação à verdade disponível dentro da família de modelo errada.

\subsection{Consequência Topológica}

Para parâmetros \textbf{contínuos} (comprimentos de ramo, taxas de substituição), convergência para o minimizador KL pode produzir valores que são ``próximos o suficiente'' para ser biologicamente úteis, mesmo que não exatamente corretos. Mas para topologia de árvore (um parâmetro \textbf{discreto}), as consequências de especificação incorreta são mais proeminentes.

Um modelo especificado incorretamente pode fazer a posterior concentrar-se em uma topologia que \textbf{não é a árvore verdadeira}, e esta concentração pode ocorrer com probabilidade posterior arbitrariamente alta conforme $n \to \infty$. Isto é o mecanismo subjacente ao fenômeno bem-documentado de inflação de probabilidade posterior: posteriores Bayesianas podem atribuir probabilidade aproximando-se a 1.0 a clados incorretos quando o modelo está errado \parencite{suzuki2002,cummings2003,yang2007}.

Em filogenômica, onde conjuntos de dados massivos produzem superfícies de verossimilhança íngremes, este modo de falha torna-se especialmente perigoso: mais dados não corrige o erro mas ao invés aumenta a certeza com qual a resposta errada é entregue \parencite{kumar2012}.

\section{Dados Finitos e \textit{a priori} Inescapável}

Embora teoremas assintóticos sejam valiosos para compreender propriedades teóricas de estimadores, \textbf{sistemática empírica lida exclusivamente com dados finitos}. O comprimento de sequência $n$ de qualquer alinhamento físico é sempre um inteiro finito.

Porque \textbf{operamos permanentemente em um regime de dados finitos}, o estado assintótico onde a verossimilhança perfeita domina \textit{a priori} nunca é estritamente atingido.

Esta equação explica este estado assintótido do qual falamos acima:

\begin{equation*}
    \lim_{t \to \infty} P_t = \pi
\end{equation*}

\textbf{Assim leríamos esta definição, passo-a-passo:}

\begin{itemize}
    \item \textbf{Expressão:} ``O limite de pê-sub-tê ($P_t$), quando o tempo tê ($t$) tende ao infinito, é igual a pi ($\pi$).''
    \item \textbf{Significado Técnico:} ``A distribuição do sistema converge para o seu estado estacionário (ou assintótico) $\pi$, independentemente de onde o sistema começou.''
    \item \textbf{Leitura Fluida:} ``Conforme o tempo passa indefinidamente, a configuração do sistema se estabiliza em uma distribuição de equilíbrio estável.''
\end{itemize}

\vspace{0.5cm}

\textbf{Algumas observações de pronúncia e contexto:}

\begin{center}
\renewcommand{\arraystretch}{1.5}
\begin{tabular}{lp{10cm}}
\hline
\textbf{Símbolo} & \textbf{Descrição e Pronúncia} \\ \hline
$\lim_{t \to \infty}$ & Lido como \textbf{``Limite com tê tendendo ao infinito''}. Representa o comportamento de longo prazo. \\
$P_t$ & A distribuição de probabilidade no instante $t$. \\
$\pi$ & Em filogenética, representa as \textbf{frequências de equilíbrio} das bases (A, C, G, T) ou aminoácidos. \\
$\to$ & A seta é lida como \textbf{``tende a''} ou \textbf{``vai para''}. \\ \hline
\end{tabular}
\end{center}

Agora, explicado o estado assintórico ao que nos referimos anteriornemte, podemos retornar à proporcionalidade do Teorema de Bayes:

$$P(\tau, \theta|X) \propto P(X|\tau, \theta) \, P(\tau, \theta)$$

\textbf{Assim leríamos esta equação, passo-a-passo:}

\begin{itemize}
    \item \textbf{Lado Esquerdo:} ``A probabilidade posterior de tau ($\tau$) e teta ($\theta$) dado xis ($X$) é proporcional a...''
    \item \textbf{Lado Direito:} ``...a verossimilhança de xis ($X$) dado tau ($\tau$) e teta ($\theta$), vezes a probabilidade a priori de tau ($\tau$) e teta ($\theta$).''
    \item \textbf{Leitura Fluida:} ``A posterior é proporcional ao produto da verossimilhança pela priori.''
\end{itemize}

\vspace{0.5cm}

\textbf{Algumas observações de pronúncia e contexto:}

\begin{center}
\renewcommand{\arraystretch}{1.5}
\begin{tabular}{lp{10cm}}
\hline
\textbf{Símbolo} & \textbf{Descrição e Pronúncia} \\ \hline
$\propto$ & Lido como \textbf{``proporcional a''}. Substitui o sinal de igual quando omitimos o denominador. \\
$\tau$ & Letra grega \textbf{``tau''}. Em filogenética, quase sempre representa a \textbf{topologia da árvore}. \\
$\theta$ & Letra grega \textbf{``teta''}. Representa os \textbf{parâmetros do modelo} (como taxas de substituição). \\
$X$ & Representa os \textbf{dados observados} (as sequências de DNA ou aminoácidos). \\ \hline
\end{tabular}
\end{center}

O posterior é sempre um compromisso matemático entre a informação contida nos dados (a verossimilhança) e a informação especificada antes dos dados serem observados (a priori). Conforme $n$ cresce, a superfície de verossimilhança torna-se mais acentuada e exerce muito mais alavancagem, aplanando o impacto relativo da priori. Porém, porque $n$ é finito, a verossimilhança nunca atinge precisão infinita.

Consequentemente, \textbf{a distribuição priori é estritamente \textit{sempre} parte da resposta final}.

\subsection{Influência Persistente da Priori}

A influência da priori (ou ``\textit{prior}'') não é uniforme através de todos os parâmetros. Para parâmetros do modelo de substituição (\textit{e.g.}, entradas da matriz de taxa e frequências de base), quantidades moderadas de dados típicamente dominam prioris razoáveis, e o efeito prático da priori torna-se negligenciável. Estes são parâmetros para os quais os dados são altamente informativos, o espaço de parâmetros é baixo-dimensional, e diferentes prioris razoáveis convergem em posteriores similares.

Porém, para outros parâmetros listados abaixo, \textit{a priori} pode persistir tenazmente:

\begin{itemize}
    \item \textbf{Topologia de árvore.} \textit{a priori} sobre topologias é uma distribuição de probabilidade sobre um espaço discreto combinatoriamente vasto. Para $n = 50$ táxons, há aproximadamente $2.8 \times 10^{76}$ árvores bifurcadas não-enraizadas. Uma priori uniforme atribui probabilidade infinitesimal a cada topologia, e diferentes prioris estruturais (uniforme vs. Yule vs. nascimento-morte) atribuem diferentes probabilidades a formas de árvore diferentes. Em regiões da árvore onde os dados carecem de sinal forte --- radiações rápidas, internós curtos, substituições saturadas --- a posterior sobre topologias concorrentes dependerá fortemente de qualquer priori que foi especificada.

    \item \textbf{Comprimentos de ramo próximos a zero.} Internós curtos, que caracterizam radiações rápidas e a ``zona de anomalia'' filogenômica, são precisamente os ramos onde os dados fornecem o mínimo de sinal e prioris tem a máxima influência. Brown \textit{et al}. (2010) demonstraram que a priori exponencial de comprimento de ramo padrão em MrBayes pode produzir probabilidades posteriores inflacionadas e estimativas de comprimento de ramo não-confiáveis, não por um bug de software, mas porque a priori era inadvertidamente informativa de formas que interagiam mal com os dados. Rannala \textit{et al}. (2012) confirmaram que prioris de comprimento de ramo podem ter efeitos substanciais em probabilidades posteriores de árvore mesmo com conjuntos de dados moderadamente grandes.

    \item \textbf{Tempos de divergência.} Em análises de relógio molecular, \textit{a priori} em taxas de ramo (\textit{e.g.}, log-normal, exponencial, ou uniforme) pode deslocar datas de divergência estimadas por dezenas de milhões de anos. Wheeler (2012) sugere que ``dada nossa falta de conhecimento, amostragem, linhagem, ambiental, e locus específicos (entre uma miríade de outros conhecidos e desconhecidos) efeitos, a distribuição uniforme pareceria tão apropriada quanto log-normal ou exponencial''. Porém, uma priori uniforme em taxas é ela própria uma afirmação forte, atribuindo probabilidade priori igual a taxas biologicamente absurdas quanto a plausíveis. Toda priori diz \textit{algo}, e ``não sei'' não é uma distribuição de probabilidade.
\end{itemize}

\section{Práticas Responsáveis para Trabalhar com Prioris e Modelos}

Se prioris sempre importam e modelos são sempre errados, a questão não é se estas limitações podem ser eliminadas (não podem). Ao invés, a questão torna-se como trabalhar responsavelmente dentro delas.

\subsection{Análise de Sensibilidade}

A resposta prática mais importante é variar ambas as prioris e modelos e examinar como os resultados mudam. Se um clado é apoiado sob uma priori mas não sob outra, ou sob um modelo mas não outro, os dados não são fortes o suficiente para apoiar aquele clado independentemente do método. Esta lógica paralela a recomendação de Wheeler (2012, Seção 10.11) para análise de sensibilidade sob regimes de custo diferente em parcimônia. Conclusões robustas a assunções analíticas são mais confiáveis que aquelas que não são.

\subsection{Teste de Adequação de Modelo}

Ao invés de assumir que o modelo escolhido é ``próximo o suficiente'', pesquisadores devem explicitamente testar se o modelo ajusta-se aos dados, usando ferramentas tais como simulações preditivas posteriores e avaliações de adequação de modelo. Se o modelo falha seu próprio teste de qualidade de ajuste, o posterior que produz é condicional em uma assunção inadequada, e nenhuma quantidade de dados a resgatará.

\begin{tcolorbox}[
    colback=red!5!white, 
    colframe=red!75!black, 
    title=\textbf{O Problema da Adequação do Modelo (\textit{Model Adequacy})},
    arc=2mm
]
\begin{itemize}
    \item \textbf{Seleção vs. Adequação:} Selecionar um modelo é como escolher o sapato que melhor serve entre três opções pequenas. Adequação é medir se o seu pé realmente cabe em algum deles. Se todos forem pequenos, o ``melhor'' ainda machucará seu pé.
    \item \textbf{Simulações Preditivas Posteriores (PPS):} É um teste de ``sanidade''. O software usa os resultados da sua análise para gerar dados sintéticos (fictícios). Se os dados simulados pelo modelo forem completamente diferentes dos seus dados reais de DNA, o modelo não ``entendeu'' o processo biológico.
    \item \textbf{O Perigo da Quantidade de Dados:} Existe um mito de que ter muitos dados (filogenômica) corrige modelos ruins. Muito pelo contrário! Se o modelo está errado, mais dados apenas darão a ele mais confiança para chegar na \textbf{resposta errada} com alta probabilidade posterior.
\end{itemize}
\end{tcolorbox}

\subsection{Relatório Honesto de Incerteza}

Uma probabilidade posterior de 1.0 para um clado significa ``este clado é esmagadoramente apoiado \textbf{dado} o modelo e as prioris.'' Não significa que ``este clado é certamente correto.'' Relatando suporte de clado juntamente com fatores de concordância que quantificam quanta proporção do genoma realmente apoia um clado dado fornecerá uma figura mais honesta da evidência que probabilidades posteriores sozinhas.

\begin{tcolorbox}[
    colback=green!5!white, 
    colframe=green!50!black, 
    title=\textbf{Posterior 1.0 vs. Realidade Biológica},
    arc=2mm
]
\begin{itemize}
    \item \textbf{A Ilusão do 1.0:} Em grandes conjuntos de dados (filogenômica), a probabilidade posterior frequentemente ``infla'' para 1.0. Isso ocorre porque o modelo Bayesiano acumula tanta evidência matemática que ele se torna ``matematicamente certo'', mesmo que o modelo em si esteja simplificando demais a biologia.
    \item \textbf{Sinal Conflitante:} O genoma não é uma unidade única. Diferentes genes podem ter histórias diferentes (devido a triagem incompleta de linhagens ou hibridização). Um posterior de 1.0 pode esconder o fato de que, por exemplo, 40\% dos dados dos seus genes apoiam uma árvore diferente.
    \item \textbf{Fatores de Concordância (CF):} Ao contrário do posterior, os CFs (como o \textit{Gene Concordance Factor}) dizem: ``De todos os genes que olhamos, quantos realmente mostram esse clado?''. Se o posterior é 1.0 mas o CF é 0.2, você sabe que o suporte é forte matematicamente, mas há um conflito biológico massivo subjacente.
\end{itemize}
\end{tcolorbox}

\section{Parcimônia Reconsiderada: Premissas para um Critério Alternativo}

As seções precedentes demonstraram que métodos probabilísticos repousam sobre assunções---especificação correta de modelo, dados infinitos, prioris bem-comportadas---que nunca são plenamente satisfeitas na prática. Isto levanta uma questão natural: pode o caso para parcimônia ser construído em seus próprios termos, não meramente como um caso especial de verossimilhança, mas como um critério mais adequado às realidades da inferência filogenética?

\subsection{Premissas}

Várias premissas sustentam este argumento. Primeiro, todos os modelos são incorretos, e incorreção do modelo não é meramente imprecisa---pode ser ativamente enganosa. Como Wheeler (2012, p.~279) afirma, a consistência de MV requer ``que o modelo usado para reconstrução seja o mesmo que gerou a árvore em primeiro lugar. Em geral, verossimilhança será inconsistente caso contrário.'' Chang (1996) mostrou que mesmo um caso tão simples quanto usar um modelo quando o cenário verdadeiro foi gerado por dois destrói a consistência \parencite{chang1996}. Segundo, consistência estatística é uma garantia assintótica que é rotineiramente violada na prática: dados empíricos de sequência não são independentes, não são identicamente distribuídos, e são finitos \parencite{wheeler2012}. Terceiro, a própria MV sofre de falhas sistemáticas: Pol e Siddall (2001) demonstraram que verossimilhança é propensa a atração de ramos longos independentemente da correção do modelo, e a ``Zona de Farris'' (Siddall, 1998)---um modo de falha não encontrado em análises de parcimônia---mostra que parcimônia pode superar verossimilhança sob condições específicas \parencite{pol2001,siddall1998}. Quarto, a história evolutiva de cada caráter é única, moldada por pressões seletivas únicas, dinâmicas populacionais e contextos mutacionais; tratar caracteres como tendo taxas evolutivas não correlacionadas (como no modelo NCM) é biologicamente mais realista do que impor uma única distribuição paramétrica compartilhada entre todos os sítios e linhagens \parencite{tuffley1997}.

\subsection{Vantagens Práticas da Parcimônia}

A partir destas premissas, várias vantagens práticas se seguem. A parcimônia é \textbf{não-paramétrica e robusta}: não assume um modelo específico de evolução de caracteres, e portanto não pode ser enganada por especificação incorreta de modelo da forma que métodos de MV e Bayesianos podem. A parcimônia \textbf{evita o problema de parâmetros de incômodo}: não requer especificar comprimentos de ramo, distribuições de taxas, ou frequências de base, todos os quais introduzem dimensões adicionais de incerteza e erro potencial. A parcimônia é \textbf{computacionalmente eficiente}: porque não requer exponenciação de matrizes ou otimização numérica de parâmetros contínuos, pontuar árvores sob parcimônia é substancialmente mais rápido, permitindo buscas mais completas do espaço de árvores. Finalmente, a parcimônia tem \textbf{poder explicativo máximo} no sentido cladístico: minimiza o número de hipóteses \textit{ad hoc} de homoplasia requeridas para explicar os dados, preferindo assim a árvore que explica mais distribuições de caracteres como eventos históricos ao invés de convergências \parencite{farris1983}.

\section{Conclusão}

Inferência filogenética probabilística é uma ferramenta matematicamente rigorosa e praticamente poderosa, mas seus resultados devem ser interpretados com uma compreensão clara de seus limites matemáticos. A convergência de posteriores à ``verdade'' é uma garantia assintótica aplicável a parâmetros contínuos condicional na assunção não-realista de especificação de modelo perfeita. Esta garantia não se estende diretamente a topologias de árvore discretas. Em realidade biológica, modelos são sempre aproximações, significando que sob especificação incorreta, inferência converge para a melhor aproximação disponível dentro da família de modelo errada. No caso da topologia de árvore, isto pode significar convergência a uma árvore incorreta com alta confiança.

Além disso, a natureza finita de quaisquer dados genômicos garante que assunções prioris continuem permanentemente um ingrediente ativo nas distribuições posteriores que calculamos, com sua influência persistindo mais fortemente precisamente onde precisamos de confiabilidade: em ramos curtos, em radiações rápidas, e em estimativas de tempo de divergência. Como nossa compreensão coletiva de métodos Bayesianos evolui com o tempo, eles podem ou não permanecer populares e considerados uma ferramenta flexível e informativa. Porém, devem ser usados com a honestidade intelectual que exigem: testando adequação de modelo, conduzindo análises de sensibilidade, e relatando incerteza de forma que distinga o que os dados dizem do que as assunções contribuem.

\clearpage
\begin{revise}\small
\begin{enumerate}[itemsep=0pt, topsep=0pt, parsep=0pt]
    \item Verossimilhança: $l(T|D) \propto \Pr(D|T)$; parâmetros de incômodo $\theta$: modelo de substituição, comprimentos de ramo, distribuição de taxas \parencite{edwards1972,wheeler2012}.
    \item Quatro formas de verossimilhança (Wheeler, 2012): \textbf{MVM} (soma sobre estados ancestrais --- o ``MV'' padrão), \textbf{MVP} (otimiza estados nos nós), \textbf{VCE} (Farris, 1973 --- especifica caminho evolutivo completo; $\text{VCE}_{\max} \equiv$ parcimônia) e \textbf{MVI} (integra sobre $\theta$; = MAP com priori plana).
    \item Paradoxo Felsenstein--Farris: resolvido reconhecendo que a ``MV'' de Felsenstein = MVM (média sobre estados ancestrais), a ``MV'' de Farris = VCE (maximização de caminhos evolutivos completos); são critérios de otimização diferentes.
    \item Modelo Sem-Mecanismo-Comum (NCM) (Tuffley \& Steel, 1997): cada caráter tem taxa única por aresta; sob $r_i$ constante, verossimilhança NCM $\equiv$ parcimônia; difere de VCE (que é variante no tratamento de estados ancestrais, não no modelo em si).
    \item Premissas para parcimônia: especificação incorreta de modelo destrói consistência de MV (Chang, 1996); MV também inconsistente na ``Zona de Farris'' (Siddall, 1998); parcimônia é não-paramétrica, evita parâmetros de incômodo, é computacionalmente eficiente e maximiza poder explicativo.
    \item Teorema de Bernstein--von Mises (BvM): posterior $\to$ normal centrada na verdade quando $n\to\infty$ --- mas \textbf{assintótico} e \textbf{condicional} ao modelo verdadeiro $\in \mathcal{M}$.
    \item Topologias são objetos discretos; espaço BHV (Billera--Holmes--Vogtmann, 2001) não é variedade suave; BvM para topologias é problema separado.
    \item Especificação incorreta ($M_{\text{verdadeiro}} \notin \mathcal{M}$): posterior converge ao minimizador KL $M^* = \arg\min D_{KL}(P_{\text{verdadeiro}} \| P_M)$, não à verdade.
    \item Consequência: mais dados $\Rightarrow$ maior certeza na resposta \textbf{errada} (inflação de probabilidade posterior) \parencite{suzuki2002,kumar2012}.
    \item Priori é \textbf{sempre} ingrediente ativo com dados finitos; influência persiste em: topologias (espaço combinatório $\sim 2{,}8\times 10^{76}$ para 50 táxons), ramos curtos (radiações rápidas) e tempos de divergência.
    \item Priori exponencial padrão de comprimento de ramo em MrBayes pode inflacionar posteriores (Brown et al., 2010; Rannala et al., 2012).
    \item Práticas responsáveis: análise de sensibilidade (variar prioris e modelos), teste de adequação de modelo (simulações preditivas posteriores), relatório honesto de incerteza (fatores de concordância).
\end{enumerate}
\end{revise}

\clearpage

\printbibliography[heading=subbibliography,title={Referências}]

